\subsection{DApp local de proves i monitorització}

La DApp local s’ha utilitzat com a interfície auxiliar per instal·lar el Snap, generar operacions de prova i verificar l’estat dels components principals del sistema. En les primeres versions del prototip, aquesta aplicació tenia una funció bàsica de connexió amb MetaMask i invocació manual del Snap. Posteriorment, es va ampliar per facilitar el seguiment de l’entorn local i actuar com a punt de control durant el desenvolupament.

En la versió actual, la DApp incorpora informació sobre l’estat del \textit{backend}, la base de dades i el servei d’Ollama, obtinguda a través d’endpoints locals com \texttt{/health} i \texttt{/stats}. Aquesta funcionalitat no forma part del mecanisme principal d’anàlisi dins MetaMask, però resulta útil per comprovar que els components necessaris estan actius abans d’executar una transacció de prova.

A més de les accions ràpides de connexió i prova del Snap, la versió actual organitza el quadre de comandament en quatre vistes: resum, historial, adreces i proves. El resum concentra l’estat dels serveis, els comptadors i les distribucions de risc i mètodes. L’historial incorpora cerca i filtre per risc; cada fila obre un detall amb el veredicte llegible i, en blocs tècnics plegables, la transacció descodificada, els indicis (\texttt{findings}), la revisió de la IA, els senyals d’adreces, les mesures de rendiment i les metadades de l’avaluació.

La mateixa interfície permet afegir manualment una adreça al conjunt local, assignar-li una etiqueta i consultar-la posteriorment. Aquesta operació utilitza els endpoints \texttt{/known-addresses/manual} i \texttt{/known-addresses}; l'historial i els detalls es recuperen mitjançant \texttt{/analysis-history} i \texttt{/analysis-history/:id}. Les etiquetes visibles de risc, acció, tipus d’adreça i procedència es mostren en català, sense alterar els valors interns de l’API. La DApp no substitueix Etherscan per generar totes les operacions contractuals: les funcions reals \texttt{approve} i \texttt{setApprovalForAll} es poden iniciar des de \textit{Write Contract}, mentre la DApp actua com a monitor local i permet activar una transacció simple per comprovar \texttt{onTransaction}. L'ús del quadre de comandament com a instrument de monitorització i les evidències visuals associades es presenten al Capítol~\ref{cap:evaluation}.
